% The q^1 pentad of phi_{0,1}^{EZ}.
% Primary anchors: Eichler-Zagier 1985, Gritsenko-Nikulin 1998,
% Gritsenko 1999, Borcherds 1995, Lorgat 2020.

\providecommand{\kBKM}{\kappa_{\mathrm{BKM}}}
\providecommand{\kch}{\kappa_{\mathrm{ch}}}
\providecommand{\kcat}{\kappa_{\mathrm{cat}}}
\providecommand{\kfib}{\kappa_{\mathrm{fiber}}}
\providecommand{\phiEZ}{\phi_{0,1}^{\mathrm{EZ}}}
\providecommand{\phiKthree}{\phi_{0,1}^{K3}}
\providecommand{\Dfive}{D_5}
\providecommand{\DeltaFive}{\Delta_5}
\providecommand{\PhiTen}{\Phi_{10}}
\providecommand{\Z}{\mathbb{Z}}
\providecommand{\C}{\mathbb{C}}

\section*{The $q^1$ Pentad of $\phi_{0,1}^{\mathrm{EZ}}$}

\paragraph{Fourier expansion.}
Let $\phiEZ(\tau,z)\in J^{\mathrm{weak}}_{0,1}$ be the
Eichler-Zagier generator in the normalization
\[
  \phiEZ(\tau,z)
  = \sum_{n\ge 0,\ \ell\in\Z}c_{\mathrm{EZ}}(4n-\ell^2)q^n r^\ell,
  \qquad q=e^{2\pi i\tau},\quad r=e^{2\pi iz}.
\]
The support condition is $4n-\ell^2\ge -1$. Through order $q^2$,
\[
\begin{aligned}
\phiEZ(\tau,z)
&=
  r^{-1}+10+r \\
&\quad
  +q(10r^{-2}-64r^{-1}+108-64r+10r^2)\\
&\quad
  +q^2(r^{-3}+108r^{-2}-513r^{-1}
       +808-513r+108r^2+r^3)
  +O(q^3).
\end{aligned}
\]
Thus
\[
 c_{\mathrm{EZ}}(-1)=1,\quad c_{\mathrm{EZ}}(0)=10,\quad
 c_{\mathrm{EZ}}(3)=-64,\quad c_{\mathrm{EZ}}(4)=108,\quad
 c_{\mathrm{EZ}}(7)=-513,\quad c_{\mathrm{EZ}}(8)=808.
\]

\paragraph{Pentad location.}
The pentad
\[
  (10,-64,108,-64,10)
\]
is the coefficient of $q^1$, read in the powers
$(r^{-2},r^{-1},r^0,r^1,r^2)$. The coefficient of $q^0$ is the triad
$(1,10,1)$, because the support condition at $n=0$ allows only
$\ell\in\{-1,0,1\}$.

\paragraph{Relation with the weight-$12$ Jacobi cusp form.}
The Fourier-Jacobi coefficient
\[
  \phi_{12,1}
  =
  (r^{-1}+10+r)q
  +(10r^{-2}-88r^{-1}-132-88r+10r^2)q^2+\cdots
\]
satisfies $\phiEZ=\phi_{12,1}/\delta_{12}$, where
\[
  \delta_{12}(\tau)
  =
  q\prod_{n\ge 1}(1-q^n)^{24}
  =
  q-24q^2+\cdots .
\]
Dividing by $\delta_{12}$ shifts the $q$-degree by one and gives
\[
\begin{aligned}
  \phiEZ|_{q^1}
  &=
  \phi_{12,1}|_{q^2}
  +24\,\phi_{12,1}|_{q^1}  \\
  &=
  (10,-88,-132,-88,10)
  +24(0,1,10,1,0)  \\
  &=
  (10,-64,108,-64,10).
\end{aligned}
\]
The weight-$12$ coefficient therefore gives the Eichler-Zagier
$q^1$ coefficient after the $\delta_{12}$ division.

\paragraph{Primitive Borcherds denominator.}
The primitive Gritsenko-Nikulin denominator is the monic product
\[
  \Dfive=64^{-1}\DeltaFive=\operatorname{Bor}(\phiEZ).
\]
Its root multiplicity convention is
\[
  \operatorname{mult}_{\mathfrak g_{\Delta_5}}(n,\ell,m)
  =
  c_{\mathrm{EZ}}(4nm-\ell^2).
\]
Borcherds' weight theorem gives
\[
  \kBKM(\DeltaFive)
  =
  \operatorname{wt}(\DeltaFive)
  =
  c_{\mathrm{EZ}}(0)/2
  =
  10/2
  =
  5.
\]
The first even null multiplicity is
$c_{\mathrm{EZ}}(0)-c_{\mathrm{EZ}}(-1)=9$.

\paragraph{K3 elliptic genus and Igusa square.}
The K3 elliptic genus is the doubled Jacobi form
\[
  \phiKthree(\tau,z)=\operatorname{Ell}(K3;\tau,z)=2\phiEZ(\tau,z).
\]
The K3-normalized $q^1$ coefficient is
$(20,-128,216,-128,20)$. The primitive denominator remains governed by
the Eichler-Zagier coefficient $(10,-64,108,-64,10)$.
Applying the same Borcherds product to the doubled input gives
\[
  \operatorname{Bor}(2\phiEZ)=\Dfive^2=4096^{-1}\DeltaFive^2,
  \qquad
  \PhiTen^{\mathrm{un}}=\DeltaFive^2,
  \qquad
  \operatorname{wt}(\PhiTen^{\mathrm{un}})=10.
\]

\paragraph{CHL and diagonal-divisor constants.}
The primitive CHL ladder is indexed by $N\in\{1,2,3,4,6\}$ and has
\[
  c_N(0)=(10,8,6,4,2),
  \qquad
  \kBKM(\Phi_N)=(5,4,3,2,1).
\]
The Gritsenko-Clery diagonal-divisor catalogue is a separate eight-row
system, with single-cusp constants
\[
  c_j^{\mathrm{GC}}(0)=(10,4,6,2,4,1,3,2),
  \qquad
  \operatorname{wt}_{\mathrm{GC}}=(5,2,3,1,2,\tfrac12,\tfrac32,1).
\]
The two row systems share the level-one $\DeltaFive$ entry.

\paragraph{$K3\times E$ invariants.}
For $X=K3\times E$, the scalar invariants are sourced separately:
\begin{center}
\begin{tabular}{c|c|l}
invariant & value & source\\
\hline
\(\kcat(X)\) & \(0\) & K\"unneth multiplicativity for \(\chi(\mathcal O_X)\)\\
\(\kch(X)\) & \(0\) & compact Hodge supertrace\\
\(\kch^{\mathrm{Heis}}(X)\) & \(3\) & Heisenberg-Mukai specialization\\
\(\kBKM(\mathfrak g_{\Delta_5})\) & \(5\) & Borcherds weight of \(\DeltaFive\)\\
\(\kfib(K3)\) & \(24\) & Mukai-lattice rank
\end{tabular}
\end{center}
These values are data of different constructions.

\paragraph{Theta-ratio computation.}
Eichler-Zagier's theta-ratio identity
\[
  \phiEZ
  =
  4\sum_{i=2,3,4}
  \left({\theta_i(\tau,z)\over\theta_i(\tau,0)}\right)^2
\]
computes the displayed coefficients. Gritsenko-Nikulin's
Fourier-Jacobi description of the Igusa product gives the same
expansion after division by $\delta_{12}$. Borcherds' denominator
theorem reads the Eichler-Zagier coefficients as root multiplicities in
the generalized Borcherds-Kac-Moody algebra attached to $\DeltaFive$.
